% Part: first-order-logic
% Chapter: syntax-and-semantics
% Section: terms-formulas

\documentclass[../../../include/open-logic-section]{subfiles}

\begin{document}

\olfileid{fol}{syn}{frm}

\olsection{Term lan \printtoken{P}{formula}}

Sawise sawijining basa logika sepisanan~$\Lang L$ diwenehi, kita bisa
nemtokake ungkapan kang kawangun saka kosakata dhasar~$\Lang L$.
Ungkapan kasebut mligine kalebu \emph{term} lan \emph{!!{formula}s}.

\begin{defn}[Term]
\ollabel{defn:terms}
Himpunan \emph{term}~$\Trm[L]$ saka~$\Lang L$
ditemtokake kanthi induktif kaya mangkene:
\begin{enumerate}
\item Saben !!{variable} iku term.
\item Saben !!{constant} saka~$\Lang L$ iku term.
\item Yen $f$ iku !!{function} $n$-panggonan lan $t_1$, \dots, $t_n$
  iku term, mula $\Atom{f}{t_1, \ldots, t_n}$ iku term.
\tagitem{limitClause}{
  Ora ana liyane kang kalebu term.}{}
\end{enumerate}
Term kang ora ngemot !!{variable}s diarani \emph{term katutup}.
\end{defn}

\begin{explain}
!!{constant}s katon minangka golongan simbol dhewe ing spesifikasi
basa lan term kita, nanging bisa uga dilebokake minangka !!{function}s
nol panggonan. Yen mangkono, kita ora mbutuhake klausa kapindho ing
definisi term. Kita mung kudu mangerteni $\Atom{f}{t_1, \ldots, t_n}$
minangka $f$ dhewe yen $n =
0$.
\end{explain}

\begin{defn}[Formula]
\ollabel{defn:formulas}
Himpunan \emph{!!{formula}s}~$\Frm[L]$ saka basa~$\Lang L$
ditemtokake kanthi induktif kaya mangkene:
\begin{enumerate}
\tagitem{prvFalse}{$\lfalse$ iku !!{formula} atomik.}{}

\tagitem{prvTrue}{$\ltrue$ iku !!{formula} atomik.}{}

\item Yen $R$ iku !!{predicate} $n$-panggonan saka~$\Lang L$ lan $t_1$, \dots,
  $t_n$ iku term saka~$\Lang L$, mula $\Atom{R}{t_1,\ldots, t_n}$ iku
  !!{formula} atomik.

\item Yen $t_1$ lan $t_2$ iku term saka~$\Lang L$, mula $\Atom{\eq}{t_1, t_2}$
  iku !!{formula} atomik.

\tagitem{prvNot}{Yen $!A$ iku !!a{formula}, mula $\lnot !A$ iku
  !!a{formula}.}{}

\tagitem{prvAnd}{Yen $!A$ lan $!B$ iku !!{formula}s, mula $(!A \land
  !B)$ iku !!a{formula}.}{}

\tagitem{prvOr}{Yen $!A$ lan $!B$ iku !!{formula}s, mula $(!A \lor !B)$
  iku !!a{formula}.}{}

\tagitem{prvIf}{Yen $!A$ lan $!B$ iku !!{formula}s, mula $(!A \lif !B)$
  iku !!a{formula}.}{}

\tagitem{prvIff}{Yen $!A$ lan $!B$ iku !!{formula}s, mula $(!A \liff !B)$
  iku !!a{formula}.}{}

\tagitem{prvAll}{Yen $!A$ iku !!a{formula} lan $x$ iku !!a{variable},
  mula $\lforall[x][!A]$ iku !!a{formula}.}{}

\tagitem{prvEx}{Yen $!A$ iku !!a{formula} lan $x$ iku !!a{variable},
  mula $\lexists[x][!A]$ iku !!a{formula}.}{}

\tagitem{limitClause}{Ora ana liyane kang kalebu !!a{formula}.}{}
\end{enumerate}
\end{defn}

\begin{explain}
Definisi himpunan term lan himpunan !!{formula}s iku
\emph{definisi induktif}. Sacara pokok, kita mbangun himpunan
!!{formula}s lumantar tahap-tahap kang cacahe tanpa wates. Ing tahap
wiwitan, kita netepake yen kabeh formula atomik iku formula; iki
cocog karo sawetara kasus wiwitan ing definisi, yaiku kasus kanggo
\iftag{prvTrue}{$\ltrue$, }{}%
\iftag{prvFalse}{$\lfalse$, }{}%
$\Atom{R}{t_1,\dots,t_n}$ lan $\Atom{\eq}{t_1,t_2}$. Dadi, ``!!{formula}
atomik'' ateges saben !!{formula} kang awujud mangkono.

Kasus liyane ing definisi menehi aturan kanggo mbangun !!{formula}s
anyar saka !!{formula}s kang wis kawangun. Ing tahap kapindho, kita
bisa nggunakake aturan kasebut kanggo mbangun !!{formula}s saka
!!{formula}s atomik. Ing tahap katelu, kita mbangun formula anyar
saka formula atomik lan formula kang kawangun ing tahap kapindho, banjur
sateruse. !!{formula} yaiku apa wae kang pungkasane kawangun ing
sawijining tahap kaya mangkono, lan ora ana liyane.
\end{explain}

Miturut konvensi, kita nulis $\eq$ ing antarane argumene lan
ngilangake kurunge: $\eq[t_1][t_2]$ iku cekakan kanggo
$\Atom{\eq}{t_1,t_2}$. Kajaba iku, $\lnot \Atom{\eq}{t_1,t_2}$
disingkat dadi $\eq/[t_1][t_2]$. Nalika nulis formula $(!B \ast !C)$
kang kawangun saka $!B$, $!C$ nganggo konektif rong panggonan~$\ast$,
kita kerep ngilangake pasangan kurung paling njaba lan mung nulis
$!B \ast !C$.

\begin{intro}
Sawetara buku teks logika mbutuhake supaya !!{variable}~$x$ kudu
muncul ing~$!A$ supaya
\iftag{prvEx}{$\lexists[x][!A]$ }{}%
\iftag{notprvEx,notprvAll}{}{lan }%
\iftag{prvAll}{$\lforall[x][!A]$ }{}%
dianggep minangka
\iftag{notprvEx,notprvAll}{!!a{formula}}{!!{formula}s}.
Yen syarat iki ora ditrapake, ora ana akibat ala, lan prekara dadi
luwih gampang.
\end{intro}

\begin{tagblock}{defNot,defOr,defAnd,defIf,defIff,defTrue,defFalse,defEx,defAll}
\begin{defn}
Formula kang kawangun nganggo operator kang ditetepake kudu dimangerteni
kaya mangkene:

\begin{tagenumerate}{defTrue,defFalse,defNot,defOr,defAnd,defIf,defIff,defEx,defAll}

\tagitem{defTrue}{$\ltrue$ nyekakake
  \iftag{prvFalse}{$\lnot\lfalse$}{$(!A \lor \lnot !A)$ kanggo sawijining
    !!{formula} atomik~$!A$ kang tetep}.}{}

\tagitem{defFalse}{$\lfalse$ nyekakake
  \iftag{prvTrue}{$\lnot\ltrue$}{$(!A \land \lnot !A)$ kanggo sawijining
    !!{formula} atomik~$!A$ kang tetep}.}{}

\tagitem{defNot}{$\lnot !A$ nyekakake $!A \lif \lfalse$.}{}

\tagitem{defOr}{$!A \lor !B$ nyekakake
  \iftag{prvAnd}{$\lnot(\lnot !A \land \lnot !B)$}{$\lnot !A \lif
    !B$}.}{}

\tagitem{defAnd}{$!A \land !B$ nyekakake
  \iftag{prvOr}{$\lnot(\lnot !A \lor \lnot !B)$}{$\lnot (!A \lif
    \lnot !B)$}.}{}

\tagitem{defIf}{$!A \lif !B$ nyekakake
  \iftag{prvOr}{$(\lnot !A \lor !B)$}{$\lnot (!A \land \lnot !B)$}.
  Cathetan koreksi sumber OLPL-064: cabang kapisan ing sumber Inggris
  ora duwe tandha kurung buka; pasangan kurunge dibalekake.}{}

\tagitem{defIff}{$!A \liff !B$ nyekakake $(!A \lif !B) \land (!B
  \lif !A)$.}{}

\tagitem{defAll}{$\lforall[x][!A]$ nyekakake $\lnot\lexists[x][\lnot !A]$.}{}

\tagitem{defEx}{$\lexists[x][!A]$ nyekakake $\lnot\lforall[x][\lnot !A]$.}{}
\end{tagenumerate}
\end{defn}
\end{tagblock}

Yen kita makarya ing basa kanggo panganggone tartamtu, kita kerep
nulis !!{predicate}s lan !!{function}s rong panggonan ing antarane
term-term argumene, umpamane $t_1 < t_2$ lan $(t_1 + t_2)$ ing basa
aritmetika sarta $t_1 \in t_2$ ing basa teori himpunan. Fungsi
penerus ing basa aritmetika malah miturut konvensi ditulis \emph{sawise}
argumene:~$t'$. Nanging, sacara resmi, iki mung cekakan konvensional
kanggo $\Atom{\Obj A^2_0}{t_1, t_2}$,
$\Obj f^2_0(t_1, t_2)$, $\Atom{\Obj A^2_0}{t_1, t_2}$ lan $\Obj f^1_0(t)$,
miturut urutane. Cathetan koreksi sumber OLPL-065: simbol fungsi
siji panggonan kang pungkasan ing sumber Inggris ora diwenehi
tipografi basa-objek kaya simbol fungsi liyane ing dhaptar iki;
tipografi kasebut dibalekake.

\begin{defn}[Identitas sintaktis]
Simbol $\ident$ mratelakake identitas sintaktis antarane rerangken
simbol, yaiku $!A \ident !B$ yen lan mung yen $!A$ lan $!B$ iku
rerangken simbol kang dawane padha lan ngemot simbol kang padha ing
saben papan.
\end{defn}

Simbol $\ident$ bisa diapit dening rerangken kang kawangun kanthi
nyambungake rerangken utawa simbol siji sawise sijine. Tuladhane,
$!A \ident (!B \lor !C)$ ateges: rerangken simbol~$!A$ padha karo
rerangken kang kawangun kanthi nyambungake, miturut urutan iki,
tandha kurung buka, rerangken $!B$, simbol $\lor$, rerangken~$!C$,
lan tandha kurung tutup. Yen mangkono, kita ngerti yen simbol kapisan
$!A$ iku tandha kurung buka, $!A$ ngemot $!B$ minangka subrerangken
(wiwit saka simbol kapindho), sawise subrerangken kasebut ana $\lor$,
lan sateruse.

Amarga term lan !!{formula}s kawangun saka unsur dhasar liwat definisi
induktif, kita bisa nggunakake prinsip-prinsip induksi iki kanggo
mbuktekake sipat-sipate.

\begin{lem}[\emph{Prinsip induksi tumrap term}]
    \ollabel{lem:trmind}
    Upamana $\Lang L$ iku basa logika sepisanan.
    Yen sawijining sipat~$P$ nyukupi syarat iki:
    %
    \begin{enumerate}
        \item sipat iku lumaku kanggo saben !!{variable}~$v$,
        %
        \item sipat iku lumaku kanggo saben !!{constant}~$a$ saka~$\Lang L$, lan
        %
        \item sipat iku lumaku kanggo $f(t_1,\dotsc,t_n)$ yen sipat iku
        lumaku kanggo $t_1$,~\dots, $t_n$ lan $f$~iku !!{function}
            $n$-panggonan saka~$\Lang L$
    \end{enumerate}
    (kanthi anggepan $t_1$,~\dots, $t_n$ iku term saka~$\Lang{L}$),
    mula $P$ lumaku kanggo saben term ing~$\Trm[L]$.
\end{lem}

\begin{prob}
    Wenehana bukti kanggo \olref[fol][syn][frm]{lem:trmind}.
\end{prob}

\begin{lem}[\emph{Prinsip induksi tumrap !!{formula}s}]
    \ollabel{thm:frmind}
    Upamana $\Lang L$ iku basa logika sepisanan.
    Yen sawijining sipat~$P$ lumaku kanggo kabeh !!{formula}s atomik
    lan nyukupi syarat iki:
    %
    \begin{enumerate}
        \tagitem{prvNot}{sipat iku lumaku kanggo $\lnot !A$ yen
            lumaku kanggo~$!A$;}{}
        \tagitem{prvAnd}{sipat iku lumaku kanggo $(!A \land !B)$
            yen lumaku kanggo $!A$ lan~$!B$;}{}
        \tagitem{prvOr}{sipat iku lumaku kanggo $(!A \lor !B)$
            yen lumaku kanggo $!A$ lan~$!B$;}{}
        \tagitem{prvIf}{sipat iku lumaku kanggo $(!A \lif !B)$
            yen lumaku kanggo $!A$ lan~$!B$;}{}
        \tagitem{prvIff}{sipat iku lumaku kanggo $(!A \liff !B)$
            yen lumaku kanggo $!A$ lan~$!B$;}{}
        \tagitem{prvEx}{sipat iku lumaku kanggo $\lexists[x][!A]$
            yen lumaku kanggo~$!A$;}{}
        \tagitem{prvAll}{sipat iku lumaku kanggo $\lforall[x][!A]$
            yen lumaku kanggo~$!A$;}{}
  \end{enumerate}
  (kanthi anggepan $!A$ lan $!B$ iku !!{formula}s saka~$\Lang{L}$),
  mula $P$ lumaku kanggo kabeh formula ing~$\Frm[L]$.
\end{lem}

\end{document}
